\clearpage
\VerdierPhysicalPageAnchor{0239}%
\begingroup\small\setlength{\parskip}{0pt}\setlength{\emergencystretch}{1.5em}
\setlength{\abovedisplayskip}{4pt plus 2pt minus 2pt}%
\setlength{\belowdisplayskip}{4pt plus 2pt minus 2pt}%

\VerdierSemanticAnchor{ch4d-4-4-2-b-cont}
\noindent\textit{このとき、}
\[
\mathrm I_2^{p,q}=H_{\mathcal B}^pF^{q,\bullet}(Y^\bullet).
\]

\VerdierSemanticAnchor{ch4d-4-4-2-c}
\noindent\textit{c) 項 $\mathrm I_\infty^{p,q}$ は、適切にフィルター付けされた
対象 $F^{p+q}(Y)$ に付随する次数付き対象である。}

\VerdierSemanticAnchor{ch4d-4-4-2-proof}
a) を証明する。$\bigl(Y(p,q),\delta(p,q,r)\bigr)$ を、複体 $Y$ に付随する、
$K(\mathcal A)$ に値をとる第一スペクトル対象とする。この対象の定義から
(4.1.2.3) により直ちに
\[
Y(-p-1,-p)=Y^p[-p].
\]
を得る。したがって（第II章 4.3.3.1）、
\[
\mathrm I_1^{p,q}=F^{p+q}\bigl(Y^p[-p]\bigr)=F^q(Y^p).
\]
微分 $d_1\colon\mathrm I_1^{p,q}\longrightarrow\mathrm I_1^{p+1,q}$ の記述は、
(4.1.2.6) から直ちに従う。主張 b) は a) から直ちに従い、最後の主張は
（第II章 4.3.2）から従う。

\VerdierSemanticAnchor{ch4d-4-4-3}
\medskip\noindent\textbf{注意 4.4.3.}\quad
$\bigl(\mathrm I_r^{p,q}(Y),d_r,r\geq1\bigr)$ を、関手 $F$ の複体 $Y$ に
関する第一スペクトル系列とする。関手
\[
Y\longmapsto\bigl(\mathrm I_r^{p,q}(Y),d_r,r\geq1\bigr)
\]
は、圏 $\operatorname{comp}(\mathcal A)$ から $\mathcal B$ に値をとる
スペクトル系列の圏への関手である。この関手は一般には圏
$K(\mathcal A)$ を通して分解しない。一方、関手

\VerdierSemanticAnchor{ch4d-4-4-3-r-ge-2}
\[
Y\longmapsto\bigl(\mathrm I_r^{p,q}(Y),d_r,r\geq2\bigr)
\]
は、命題 (4.4.2) b) により、複体のホモトピー圏を通して分解する。

\VerdierSemanticAnchor{ch4d-4-4-4}
\medskip\noindent\textbf{4.4.4.}\quad
今度は $\mathcal A$ をアーベル圏、
$F\colon D(\mathcal A)\longrightarrow\mathcal B$ をアーベル圏
$\mathcal B$ に値をとるコホモロジー的関手とし、$F^p$、$p\in\mathbb Z$ を、
$F$ と $D(\mathcal A)$ のシフトとの合成から得られる関手
（$F^p(X)=F(X[p])$）とする。さらに $Y$ を
$\operatorname{comp}(\mathcal A)$ の対象とし、
$\bigl(Y(p,q),\delta(p,q,r)\bigr)$ を、複体 $Y$ に付随する
$D(\mathcal A)$ に値をとる第一スペクトル対象とする。関手 $F$ はこの
スペクトル対象を $\mathcal B$ に値をとるスペクトル対象に移す。これから
$\mathcal B$ に値をとるスペクトル系列を得るが、これも
\emph{関手 $F$ の複体 $Y$ に関する第一スペクトル系列}と呼ぶ。この
スペクトル系列は本質的に命題 (4.4.2) によって記述される（命題 (4.4.2) で
圏 $K(\mathcal A)$ を圏 $D(\mathcal A)$ に置き換えればよい）。関手 $F$ の
複体 $Y$ に関する第一スペクトル系列は、$F$ について、また論元 $Y$ が
$\operatorname{comp}(\mathcal A)$ の中を動くときのその論元について関手的である。
このスペクトル系列には常に（第II章 4.3.3.1）で導入した添字付けを用いる。
\endgroup
